\documentclass[UTF8,a4paper,12pt]{ctexart}
\usepackage{iftex}
\usepackage{geometry}
\usepackage{setspace}
\usepackage{amsmath,amssymb,booktabs,graphicx,float}
\usepackage{caption,subcaption}
\usepackage{enumitem}
\usepackage{longtable}
\usepackage{array}
\usepackage{xcolor}
\usepackage{hyperref}

\geometry{left=2.5cm,right=2.5cm,top=2.5cm,bottom=2.5cm}
\setstretch{1.25}
\setlength{\parindent}{2em}
\setlength{\parskip}{0pt}
\graphicspath{{../outputs/figures/}{figures/}{../figures/}{outputs/figures/}}

\ifXeTeX
  \setmainfont{Times New Roman}
  \IfFontExistsTF{SimSun}{\setCJKmainfont{SimSun}[AutoFakeBold=2.5]}{}
  \IfFontExistsTF{SimHei}{\setCJKsansfont{SimHei}}{}
  \IfFontExistsTF{FangSong}{\setCJKmonofont{FangSong}}{}
\fi

\hypersetup{
  hidelinks,
  pdfborder={0 0 0},
  pdftitle={论文题目},
  pdfauthor={},
  pdfsubject={CUMCM mathematical modeling paper}
}

\title{\heiti 论文题目}
\author{}
\date{}

\begin{document}
\pagenumbering{arabic}
\maketitle
\thispagestyle{plain}

\begin{abstract}
这里写一页内摘要。按问题顺序说明模型、方法、关键结果和结论。所有数值结论必须来自已确认结果，并能在 \texttt{outputs/result\_contract.json} 中追溯到脚本、参数和输出文件。

\textbf{关键词：} 关键词一；关键词二；关键词三；数学建模
\end{abstract}
\newpage

\section{问题重述}
用自己的语言重述任务、输入、输出、约束和评价目标，不照抄题面。

\section{问题分析}
说明每个子问题的求解逻辑和模块关系。

\begin{figure}[H]
  \centering
  \includegraphics[width=0.86\textwidth]{modeling_workflow.pdf}
  \caption{整体建模流程图。绘图完成后只需替换文件名或放入同名文件。}
  \label{fig:workflow}
\end{figure}

\section{模型假设}
\begin{enumerate}[label=H\arabic*:]
  \item 假设一：说明必要性、合理性和可能影响。
  \item 假设二：说明必要性、合理性和可能影响。
\end{enumerate}

\section{符号说明}
\begin{table}[H]
  \centering
  \caption{主要符号说明}
  \label{tab:notation}
  \begin{tabular}{lll}
    \toprule
    符号 & 含义 & 单位 \\
    \midrule
    $x_i$ & 待填写变量含义 & 待填写 \\
    $y_i$ & 待填写变量含义 & 待填写 \\
    \bottomrule
  \end{tabular}
\end{table}

\section{数据处理与探索}
说明数据来源、缺失值、异常值、标准化和特征构造。数据来源、许可和处理脚本需在支撑材料中可追溯。

\begin{table}[H]
  \centering
  \caption{数据概览与处理说明}
  \label{tab:data-overview}
  \begin{tabular}{llll}
    \toprule
    变量 & 单位 & 问题 & 处理方式 \\
    \midrule
    待填写 & 待填写 & 缺失/异常/尺度 & 待填写 \\
    \bottomrule
  \end{tabular}
\end{table}

\section{模型建立与求解}

\subsection{问题一模型}
写明目标函数、约束、算法步骤和输出。

\begin{equation}
  \min f(x) = \sum_i c_i x_i
\end{equation}

\begin{figure}[H]
  \centering
  \includegraphics[width=0.82\textwidth]{q1_result.png}
  \caption{问题一结果图。待代码生成后替换。}
  \label{fig:q1-result}
\end{figure}

\subsection{问题二模型}
写明模型结构、参数估计、评价指标和求解流程。

\begin{figure}[H]
  \centering
  \includegraphics[width=0.82\textwidth]{q2_result.png}
  \caption{问题二结果图。待代码生成后替换。}
  \label{fig:q2-result}
\end{figure}

\subsection{问题三模型}
写明拓展模型、优化方案或综合评价。

\section{结果分析与检验}
\begin{table}[H]
  \centering
  \caption{核心结果汇总}
  \label{tab:result-summary}
  \begin{tabular}{llll}
    \toprule
    子问题 & 指标 & 结果值 & 结论 \\
    \midrule
    问题一 & 待填写 & 待确认 & 待确认 \\
    问题二 & 待填写 & 待确认 & 待确认 \\
    问题三 & 待填写 & 待确认 & 待确认 \\
    \bottomrule
  \end{tabular}
\end{table}

\section{灵敏度、误差与稳健性分析}
\begin{figure}[H]
  \centering
  \includegraphics[width=0.82\textwidth]{sensitivity.png}
  \caption{关键参数灵敏度分析图。待代码生成后替换。}
  \label{fig:sensitivity}
\end{figure}

\section{模型评价、改进与推广}
客观说明优点、局限、适用范围和后续改进方向。

\begin{thebibliography}{99}
\bibitem{ref1} 作者. 文献题名. 来源, 年份.
\end{thebibliography}

\appendix
\section{支撑材料说明}
列出代码文件、数据文件、图表文件、AI 工具使用说明和三个机器可读合同文件：\texttt{result\_contract.json}、\texttt{constraint\_checks.json}、\texttt{figure\_manifest.json}。完整代码放入支撑材料。

\end{document}
